\documentclass[10pt,a4paper]{amsart}
\usepackage[margin=19mm]{geometry}
\usepackage{amsmath,amssymb,mathtools}
\usepackage[scr=rsfs,cal=euler]{mathalfa}
\usepackage{xfrac}
\usepackage[unicode,hidelinks,bookmarks=false]{hyperref}
\newcommand{\ie}{i.e.\ }
\newcommand{\cf}{cf.\ }
\newcommand{\resp}{respectively\ }
\newcommand{\Subsection}{\subsection}
\providecommand{\nobreakdash}{\nobreak\hbox{-}}
\setlength{\emergencystretch}{2em}
\multlinegap0pt
\usepackage{bm}
\usepackage{bbm}
\usepackage{dsfont}
\usepackage{xspace}
\usepackage{cancel}
\usepackage{mathtools}
\usepackage{amsthm}
\makeatletter
\@ifundefined{lemma}{\newtheorem{lemma}{Lemma}}{}
\@ifundefined{remark}{\newtheorem{remark}{Remark}}{}
\makeatother
\title[Explicit Formulas For Generalized Polylogarithmic Integrals,]{Explicit Formulas For Generalized Polylogarithmic Integrals, Euler Sums, And BBP-Type Series}
\author{Ali Olaikhan}
\date{Source: arXiv:2507.04205v1 (CC BY 4.0)}
\begin{document}
\maketitle

\noindent\emph{Source and licence.} Explicit Formulas For Generalized Polylogarithmic Integrals, Euler Sums, And BBP-Type Series. Version arXiv:2507.04205v1 (CC BY 4.0), distributed under the Creative Commons Attribution 4.0 International licence, as recorded in the arXiv licence metadata for this exact version and shown on the abstract page https://arxiv.org/abs/2507.04205v1. Full rights notice: licenses/arxiv-2507.04205-CC-BY-4.0.txt.\par\smallskip

\noindent\textit{Editorial scope.} Counted unit: the Remark whose source label is \texttt{None} in the article (source0.tex lines 309--311, source label \texttt{None}), delivered together with the article's own complete original proof of it (source0.tex lines 312--342, one \texttt{proof} environment of 31 lines). No mathematics is written, repaired or re-derived: statement and proof are the article's own text line for line. Required context retained uncounted: polylog (lines 83--85); lerch integral (lines 210--213); theta integral (lines 226--235); tough integral (lines 246--251). Every definition, display and label of those lines is retained unchanged. Omitted from this package and not redistributed: 1--82, 86--209, 214--225, 236--245, 252--308, 343--826. Nothing inside a retained excerpt is omitted or altered: every retained body is delivered exactly as the article's lines are written, so no omission receipt of any kind accompanies this package. Citations: the retained excerpts carry no citation command at all, so no bibliography entry of the article is transcribed and no reference is redistributed. Reference adaptation: none was needed and none was made; every reference inside the delivered unit resolves inside it, so no unresolved reference or citation is produced. Printed slips kept literal: no printed word, symbol, hypothesis or number of the counted statement or of its proof was altered. Layout and class: the article's own class is \texttt{article} with packages including color, amssymb, amsmath, amsthm, amsfonts, amscd, graphicx, hyperref, fullpage, float, psfig, graphics, latexsym, epsf; the wrapper is the lane's pinned \texttt{amsart} editorial wrapper at 10pt on A4, which restates the theorem-like environments the retained text needs and adds the article's own macro definitions that the retained text needs (none), with extra wrapper packages (bm, bbm, dsfont, xspace, cancel, mathtools). The delivered numbering is the unit's own. Assets: no retained span contains a figure environment, a graphics-inclusion command, a PDF-page inclusion, a table or a TikZ picture; no image file is copied into this package, the package declares no asset dependency and no image file is redistributed with it. Licence: version arXiv:2507.04205v1 of this article is distributed under the Creative Commons Attribution 4.0 International licence, as recorded in the arXiv licence metadata for this exact version and shown on the abstract page https://arxiv.org/abs/2507.04205v1. A full rights notice is delivered at licenses/arxiv-2507.04205-CC-BY-4.0.txt. Characters: every retained excerpt is pure ASCII, so the delivered engine, XeLaTeX with its default Latin Modern text font, reports no missing character.
\begin{equation}
	\operatorname{Li}_p(x)=\frac{{(-1)}^{p-1}}{(p-1)!}\int_0^1\frac{x\ln^{p-1}(y)}{1-xy}\mathrm{d} y,\quad x\leqslant1.\label{polylog}
\end{equation}

\begin{lemma}\label{lerch integral}
	Let $q\in\mathbb{Z}^{+}$, $s,r>0$, and $c=\pm1$, with $q\neq1$ when $c=1$. Then
	\[\int_0^1\frac{x^{s-1}\ln^{q-1}(x)}{1-cx^r}\mathrm{d} x=\frac{{(-1)}^{q-1}(q-1)!}{r^{q}}\,\Phi\left(c,q,\tfrac{s}{r}\right).\]
\end{lemma}

\begin{lemma} \label{theta integral}
	Let $q\in\mathbb{Z}^+$, $c=\pm1$, and $\frac{s}{r}\in(0,1)$. Then
	\[\int_0^\infty\frac{x^{s-1}\ln^{q-1}(x)}{1-cx^r}\mathrm{d} x=\frac{{(-1)}^{q-1}(q-1)!}{r^{q}}\,\Theta(c,q,s,r),\]
	where
	\begin{equation}
		\Theta(c,q,s,r)=\Phi\left(c,q,\tfrac{s}{r}\right)+c{(-1)}^{q}\Phi\left(c,q,\tfrac{r-s}{r}\right),\label{theta}
	\end{equation}
	from which it follows that
	\[\Theta(c,q,1,2)=(1+c{(-1)}^{q})\,\Phi(c,q,\tfrac12).\]
\end{lemma}

\begin{lemma} \label{tough integral}
	Let $q,n\in\mathbb{Z}^{+}$, $y\in(0,1)$, $a,b=\pm1$, with $q\neq1$ when $a=1$. Then
	\[\int_0^\infty \frac{\ln^{q-1}(x)}{x}\left(\frac{1}{1-ax}-\frac{1}{1-byx^{n}}\right)\mathrm{d} x\]
	\[=-a(1+{(-1)}^{q})(q-1)!\,\Phi(a,q,1)
	+\frac{2b{(-1)}^{q} (q-1)!}{n^{q}}\sum_{j=0}^{\lfloor\frac{q}{2}\rfloor}\frac{\ln^{q-2j}(y)}{(q-2j)!}\,\Phi(b,2j,1).\]
\end{lemma}

\begin{remark}
	The case $p=a^nb=1$ is valid only when $n=1$, which leads to $a=b$, as the first term containing the factor $\Phi(a^nb,p,1)$ cancels with the last term in the first summation. The integral must be understood as the Cauchy principal value when $b=1$.
\end{remark}

\begin{proof}
	We have in (\ref{polylog}) that $x\leqslant1$. However, the polylogarithm can be extended to $x>1$ if the Cauchy principal value of the integral is taken into consideration, that is
	\begin{equation}
		\operatorname{Li}_p(x) = \frac{{(-1)}^{p-1}}{(p-1)!} \,\operatorname{P.V.} \int_0^1 \frac{x \ln^{p-1}(y)}{1 - xy} \, \mathrm{d} y, \quad x > 1.\label{pv}
	\end{equation}
	Using this definition, we write
	\[\int_0^\infty\frac{\ln^{q-1}(x)\operatorname{Li}_p(bx^n)}{x(1-ax)}\mathrm{d} x=\frac{{(-1)}^{p-1}}{(p-1)!}\int_0^\infty\int_0^1\frac{bx^n\ln^{q-1}(x)\ln^{p-1}(y)}{x(1-ax)(1-byx^n)}\mathrm{d} y\mathrm{d} x,\]
	where the double integral must be understood as the Cauchy principal value for the case $m=1$, as stated in (\ref{pv}). 
	Next, we swap the order of integration, which is justified by Fubini's theorem. Therefore, the main integral becomes
	\begin{equation}
		\int_0^\infty\frac{\ln^{q-1}(x)\operatorname{Li}_p(bx^n)}{x(1-ax)}\mathrm{d} x=\frac{{(-1)}^{p-1}}{(p-1)!}\int_0^1\ln^{p-1}(y)\int_0^\infty\frac{bx^n\ln^{q-1}(x)}{x(1-ax)(1-byx^n)}\mathrm{d} x\mathrm{d} y.\label{main}
	\end{equation}
	By partial fraction decomposition we have for $a,b=\pm1$:
	\begin{equation}		\frac{x^{n}}{(1-ax)(1-byx^{n})}=\frac{a^n}{1-a^n by}\left(\frac{1}{1-ax}-\frac{\sum_{j=1}^n (ax)^{j-1}}{1-byx^{n}}\right),\label{pattern}
	\end{equation}
	from which it follows that
	\[\int_0^\infty\frac{x^{n}\ln^{q-1}(x)}{x(1-ax)(1-byx^{n})}\mathrm{d} x=\frac{a^n}{1-a^n by}\int_0^\infty\frac{\ln^{q-1}(x)}{x}\left(\frac{1}{1-ax}-\frac{\sum_{j=1}^n (ax)^{j-1}}{1-byx^{n}}\right)\mathrm{d} x.\]
	Breaking up the integrand does solve the problem, as both resulting integrals would diverge. Instead, we separate the first term of the inner sum and then change the order of integration and summation. This yields
	\[\int_0^\infty\frac{x^{n}\ln^{q-1}(x)}{x(1-ax)(1-byx^{n})}\mathrm{d} x=\frac{a^n}{1-a^n by}\int_0^\infty\frac{\ln^{q-1}(x)}{x}\left(\frac{1}{1-ax}-\frac{1}{1-byx^n}\right)\mathrm{d} x\]
	\[-\frac{a^n}{1-a^n by}\sum_{j=2}^n a^{j-1}\int_0^\infty\frac{x^{j-2}\ln^{q-1}(x)}{1-byx^n}\mathrm{d} x.\]
	With these adjustments, both integrals are now convergent. The first integral appears in Lemma \ref{tough integral}. For the second integral, we let $yx^n=t^n$, apply the binomial expansion for $\ln^{q-1}(y^{-\frac{1}{n}}t)$, and then make use of Lemma \ref{theta integral}, we find
	\begin{equation}
		\int_0^\infty\frac{x^{j-2}\ln^{q-1}(x)}{1-byx^{n}}\mathrm{d} x=\frac{{(-1)}^{q-1}(q-1)!}{n^{q}}\sum_{k=1}^{q}\frac{\Theta(b,k,j-1,n)}{(q-k)!}y^{\frac{-j+1}{n}}\ln^{q-k}(y).\label{first integral}
	\end{equation}
	Collecting these two integrals, we conclude
	\[\int_0^\infty\frac{ x^{n-1}\ln^{q-1}(x)}{(1-ax)(1-byx^{n})}\mathrm{d} x
	=-\frac{a^{n-1}(q-1)!}{1-a^n by}(1+{(-1)}^{q})\,\Phi(a,q,1)\]
	\[+\frac{2a^nb{(-1)}^{q} (q-1)!}{(1-a^n by)n^{q}}\sum_{j=0}^{\lfloor\frac{q}{2}\rfloor}\frac{\Phi(b,2j,1)}{(q-2j)!}\ln^{q-2j}(y)\]
	\[+\frac{a^{n-1}{(-1)}^{q}(q-1)!}{(1-a^n by)n^{q}}\sum_{j=2}^n \sum_{k=1}^{q}a^{j}\frac{\Theta(b,k,j-1,n)}{(q-k)!}y^{\frac{-j+1}{n}}\ln^{q-k}(y).\]
	Returning to (\ref{main}) and then using Lemma \ref{lerch integral} completes the proof.
\end{proof}

\end{document}
